\chapter{复线性空间}
\label{chap:complex linear space}

这一章我们将数域由实数域$\RR$扩展至复数域$\CC$，复数的定义和运算高中已经讲过，也可参见复变函数的笔记。在此略。

复数构成的向量$z$的共轭即将其中所有元素取共轭，记作$\bar z$。共轭转置记作$z\dg:=\bar z\tp$。

线性无关、维数和线性映射等代数概念可推广到复数域；内积、正交性和谱定理则必须加入复共轭，不能仅将实数符号机械地替换为复数。以下涉及基、对偶和谱分解的结论均在有限维空间中讨论。
\section{内积和内积空间}
\begin{definition}{$\CC^n$标准内积}{standard inner product}
	复向量$u,v$的内积
	\[
		u\dg v=\sum_{i=1}^n\bar u_iv_i=\bar u_1v_1+\cdots+\bar u_nv_n.
	\]
	\tcblower
	一般复线性空间$V$的内积$\inp\cdot\cdot:V\times V\to\CC$：
	\begin{compactitem}
		\item 交换共轭：$\inp uv=\overline{\inp vu};$
		\item 对第二个变量线性：$\inp u{cv}=c\inp uv,\enspace \inp u{v+w}=\inp uv+\inp uw;$
		\item 正定：$\inp uu\geqslant 0$当且仅当$u=0$时取等号。
	\end{compactitem}
\end{definition}
\begin{remark}
	对第一个变量不是简单的线性，而是多一个复共轭：
	\[
		\inp{cu}v=\bar c\inp uv,\quad \inp{u+v}w=\inp uw+\inp vw.
	\]
\end{remark}
% 定义了内积的空间叫做内积空间(inner product space)。
\begin{theorem}{}{}
	只需要知道基之间的内积就可以算出任意向量之间的内积。
\end{theorem}
\begin{proof}
	$v_1,\ldots,v_n$是$V$上的一组基，%$g_{ij}:=\inp{v_i}{v_j}$，则$g_{ji}=\bar g_{ij}.$
	$\forall u,w\in V,\enspace u=u^1v_1+\cdots+u^nv_n,\enspace w=w^1v_1+\cdots+w^nv_n$
	\[
		\inp uw=\inp{u^1v_1+\cdots+u^nv_n}{w^1v_1+\cdots+w^nv_n}=\sum_{i,j}\bar u^iw^j\inp{v_i}{v_j}.
		\qedhere
	\]
\end{proof}
\begin{example}{内积与对偶空间}{inner product and dual space}
	$V$的对偶空间$V^\ast$是所有$V\to\CC$的线性函数的集合。通过内积可以建立$V,V^\ast$的一一映射（共轭线性，而非复线性）
	\[
		\forall v\in V,g_v\in V^\ast,\enspace g_v(w):=\inp vw
.	\]
\end{example}
\begin{example}{Legendre多项式}{Legendre polynomial}
	所有不高于$n$的实系数多项式
	\[
		f(x)=a_0+a_1x+\cdots+a_nx^n,
	\]
	构成线性空间$\mathscr P^n(\RR)$，显然$\{1,x,x^2,\ldots,x^n\}$构成$\mathscr P^n(\RR)$的一组基。定义内积
	\[
		\inp fg:=\int_{-1}^1f(x)g(x)\d x,
	\]
	用Gram-Schmidt法则将$\{1,x,x^2,\ldots,x^n\}$变成一组正交基%$\{P_0,P_1,P_2,\ldots,P_n\}$
	\begin{align*}
		P_0&=1,\\
		P_1&=x-\frac{\inp{P_0}x}{\inp{P_0}{P_0}}P_0=x,\\
		P_2&=x^2-\frac{\inp{P_1}{x^2}}{\inp{P_1}{P_1}}P_1-\frac{\inp{P_0}{x^2}}{\inp{P_0}{P_0}}P_0=x^2-\frac13,\\
		P_3&=x^3-\frac{\inp{P_2}{x^3}}{\inp{P_2}{P_2}}P_2-\frac{\inp{P_1}{x^3}}{\inp{P_1}{P_1}}P_1-\frac{\inp{P_0}{x^3}}{\inp{P_0}{P_0}}P_0=x^3-\frac35x,\\
		&\cdots
	\end{align*}
	这里得到首一多项式；标准Legendre多项式通常取$P_n(1)=1$，例如标准$P_2(x)=(3x^2-1)/2$，与上述首一版本只差归一化系数。
\end{example}
\begin{example}{Hermite多项式}{Hermite polynomial}
	在$\mathscr P^n(\RR)$内定义内积
	\[
		\inp fg:=\int\iti f(x)g(x)\e^{-x^2/2}\d x,
	\]
	用Gram-Schmidt法则将$\{1,x,x^2,\ldots,x^n\}$变成一组正交基
	\begin{align*}
        \mathrm{He}_0&=1, &\mathrm{He}_1&=x,\\
        \mathrm{He}_2&=x^2-1, &\mathrm{He}_3&=x^3-3x.
	\end{align*}
	这里得到的是概率论约定的Hermite多项式$\mathrm{He}_n$。物理学中常用的$H_n$对应权重$\e^{-x^2}$，二者还涉及自变量缩放：$H_n(x)=2^{n/2}\mathrm{He}_n(\sqrt2x)$，不能仅相差一个整体系数。
\end{example}

\section{幺正矩阵}

在涉及内积的实矩阵公式中，应将转置换为共轭转置。例如复矩阵的SVD为$A=U\Sigma V\dg$，其奇异值由$A\dg A$决定。

\begin{definition}{幺正矩阵}{unitary matrix}
	方阵$U\in\CC^{n\times n}$是幺正的(unitary)若$U\dg U=I$；此时也有$UU\dg=I$。矩形矩阵仅满足$U\dg U=I$时是等距嵌入，不称为幺正矩阵。
\end{definition}
幺正矩阵是正交矩阵(见\dfnref{dfn:orthogonal matrix})在复空间的推广。
\begin{theorem}{幺正变换}{unitary transfomation}
	幺正变换保持复向量的模不变。
\end{theorem}
\begin{proof}
	\[
		\norm{Uz}^2=z\dg U\dg Uz=z\dg z=\norm z^2.
		\qedhere
	\]
\end{proof}
\begin{theorem}{幺正矩阵的行列式}{determinant of unitary}
	$\abs{\det(U)}=1.$
\end{theorem}
\begin{proof}
	\[
		1=\det(U\dg U)=\det(U\dg)\det(U)=\overline{\det(U)}\det(U)=\abs{\det(U)}^2.
		\qedhere
	\]
\end{proof}

\section{Hermite矩阵}

\begin{definition}{Hermite矩阵}{Hermite matrix}
	方阵$H$是厄米(Hermite)矩阵若$H\dg=H.$
\end{definition}
Hermite矩阵是实对称矩阵(见\dfnref{dfn:symmetric matrix})在复空间的推广。

\begin{example}{Pauli矩阵}{Pauli matrix}
	给出三个Pauli矩阵
	\[
		\sigma_1=
		\begin{bmatrix}
			0&1\\1&0
		\end{bmatrix},\enspace
		\sigma_2=
		\begin{bmatrix}
			0&-\i\\\i&0
		\end{bmatrix},\enspace
		\sigma_3=
		\begin{bmatrix}
			1&0\\0&-1
		\end{bmatrix}.
	\]
	$\sigma_1,\sigma_2,\sigma_3$都是Hermite的，且
	\[
		\sigma_i\sigma_j=\i\sigma_k,\quad(ijk)=(123).
	\]
\end{example}

\begin{theorem}{Hermite矩阵的二次型}{quadratic form of Hermite}
	$\forall z\in\CC^n$，$z\dg Hz$是实数。
\end{theorem}
\begin{proof}
	$(z\dg Hz)\dg=z\dg H\dg z=z\dg Hz.$
\end{proof}
\begin{theorem}{Hermite矩阵的特征值}{}
	Hermite矩阵$H$的特征值都是实数。
\end{theorem}
\begin{proof}
	$Hz=\lambda z$，左乘$z\dg$得$z\dg Hz=\lambda z\dg z$，由$z\dg Hz,z\dg z$均是实数知，$\lambda$也是实数。
\end{proof}
\begin{theorem}{Hermite矩阵的特征向量}{eigenvalue of Hermite}
	Hermite矩阵$H$不同特征值对应的特征向量正交。
\end{theorem}
\begin{proof}
	$Hz_1=\lambda_1z_1,\enspace Hz_2=\lambda_2z_2,\quad\lambda_1\neq\lambda_2$
	\[
		\lambda_1z_2\dg z_1=z_2\dg Hz_1=(z_1\dg Hz_2)\dg=(\lambda_2z_1\dg z_2)\dg=\lambda_2z_2\dg z_1.
	\]
	故$z_2\dg z_1=0.$
\end{proof}
\begin{theorem}{谱定理}{spectral theorem of Hermite}
	Hermite矩阵可以选出$n$个特征向量$u_1,\ldots,u_n$构成$\CC^n$的一组正交归一基。重复特征值的子空间内需要另作正交归一化，并非任意选取的特征向量都正交。令$U=[u_1,\ldots,u_n]$，则
	\[
		H=U\Lambda U\dg.
	\]
\end{theorem}
\begin{proof}
    由代数基本定理，$H$有归一化特征向量$u_1$，特征值$\lambda_1$为实数。若$w\perp u_1$，则
    \[
        \inp{u_1}{Hw}=\inp{Hu_1}{w}=\lambda_1\inp{u_1}{w}=0.
    \]
    因而$u_1^\perp$在$H$作用下不变，$H$在该子空间上的限制在正交归一基下仍有Hermite矩阵表示。对维数归纳，即可在整个空间选出正交归一特征基。
\end{proof}

\section{伴随、自伴与正规映射}
\begin{definition}{伴随映射}{adjoint map}
    有限维复内积空间上的线性映射$T:V\to V$的伴随$T^\dagger$由
    \[
        \inp{Tu}{v}=\inp{u}{T^\dagger v},\qquad u,v\in V
    \]
    唯一确定。若$T=T^\dagger$，则称$T$自伴。
\end{definition}
在正交归一基中，伴随映射的矩阵恰为原矩阵的共轭转置。因此有限维自伴映射与Hermite矩阵对应。若所选基并不正交归一，记Gram矩阵$G_{ij}=\inp{v_i}{v_j}$，$T$的矩阵为$A$，则伴随的矩阵为
\[
    [T^\dagger]=G^{-1}A^\dagger G.
\]
此时自伴条件是$A^\dagger G=GA$，不能直接要求$A^\dagger=A$。无限维空间的无界算符还需要指定定义域，不能直接照搬这里的有限维判定。

\begin{definition}{正规矩阵}{normal matrix}
    方阵$A\in\CC^{n\times n}$满足$A^\dagger A=AA^\dagger$时，称为正规矩阵。Hermite矩阵和幺正矩阵都是正规矩阵。
\end{definition}
\begin{theorem}{正规矩阵的谱定理}{spectral theorem of normal matrix}
    复方阵$A$能被幺正矩阵对角化，当且仅当$A$正规。
\end{theorem}
\begin{proof}
    若$A=U\Lambda U^\dagger$且$\Lambda$对角，直接计算可得$A^\dagger A=AA^\dagger$。
    反之，取$A$的归一化特征向量$q_1$，将其补成正交归一基，得幺正矩阵$Q$与
    \[
        Q^\dagger AQ=\begin{bmatrix}\lambda&b^\dagger\\0&B\end{bmatrix}=:C.
    \]
    幺正相似变换保持正规性。比较$CC^\dagger=C^\dagger C$的左上元素，得$|\lambda|^2+\norm b^2=|\lambda|^2$，故$b=0$。此时$B$也正规，对维数归纳即可完成幺正对角化。
\end{proof}
\begin{remark}
    一般可对角化矩阵未必正规。例如$\begin{bmatrix}1&1\\0&2\end{bmatrix}$有两个不同特征值，因而可相似对角化，但不能幺正对角化。实正交矩阵也未必能在实数域上对角化，例如平面上转角$\pi/2$的旋转矩阵；其复特征值为$\pm\i$。
\end{remark}
